\section{A quantitative regularity budget for stationary actions}
\label{sec:v80-regularity}

We make explicit the finite-order hypotheses for the convex-twist
inverse.  All configuration and phase domains in this section are the
fixed compact domains constructed in the proof of
Theorem~\ref{thm:v79-six-clock}, enlarged by a fixed positive collar.
The number of iterates is fixed.  Constants do not depend on the
particular generating function within the declared class.

\begin{proposition}[Propagation of finite-order bounds]
\label{prop:v80-stationary-bounds}
Fix $q\ge1$, $m\ge0$ and the convexity constants in
\eqref{eq:v79-convex-class}.  Let the minimizing paths for endpoints in
a fixed compact gate have all coordinates in $[-K,K]$.  Suppose
\[
              \|L\|_{C^{m+2}([-K-1,K+1]^2)}\le M_L.
\]
Then the interior minimizing coordinates $X_q$ have a uniform
$C^{m+1}$ bound, and the stationary action $S_q$ has a uniform
$C^{m+2}$ bound on the gate.  The constants depend only on
$q,m,K,M_L,\lambda$ and the gate.  For raw source and efficiency bounds
\[
 \|\rho_q\|_{C^m}\le M_\rho,\qquad \|e_q\|_{C^m}\le M_e
\]
on their fixed larger domains, the weight
\begin{equation}\label{eq:v80-weight-budget}
 a_q=D^{-1}e_q(s,t)\rho_q(s,-(S_q)_s(s,t))|(S_q)_{st}(s,t)|
\end{equation}
and the raw clock subdensities have uniform $C^m$ bounds.  Positive
source and efficiency floors give a uniform positive weight floor
using~\eqref{eq:v79-twist-bound}.  Conditional densities have the same
regularity when their normalizing masses have a positive lower bound.
\end{proposition}
\begin{proof}
The case $q=1$ has no interior variables and follows directly.
For $q>1$ set
\[
 \mathcal A(z,X)=\sum_{i=0}^{q-1}L(x_i,x_{i+1}),\quad
 z=(s,t),\quad G=\partial_X\mathcal A.
\]
The stationarity equation is $G(z,X_q(z))=0$, and its $X$ derivative
$H$ satisfies $H\succeq2\lambda I$.  Hence
$\|H^{-1}\|\le(2\lambda)^{-1}$.  Derivatives of $G$ through order
$m+1$ are bounded by constants times $qM_L$.

Here is an explicit recursive form of the implicit bound.  Put
$P_0=\sqrt{q-1}K$.  At first order use
$D X_q=-H^{-1}G_z$.  At order $j\ge2$, differentiation of the
stationarity equation gives
\[
 H D^jX_q=-\mathcal B_j,
\]
where $\mathcal B_j$ is the finite chain-rule sum indexed by partitions
of the $j$ differentiations, after removing the single term
$H D^jX_q$.  Every remaining term contains a derivative of $G$ of
order at most $j$ and derivatives of $X_q$ of orders strictly less than
$j$.  Let $c_{q,j}$ be the sum of the absolute integer coefficients of
that finite multivariable chain-rule sum, enlarged if necessary to cover
its coordinate norms and the first-order formula.  The recursive choice
\begin{equation}\label{eq:v80-implicit-recursion}
 P_j=\frac{c_{q,j}qM_L}{2\lambda}
                   \left(1+\sum_{i=1}^{j-1}P_i\right)^j,
                    \qquad 1\le j\le m+1,
\end{equation}
bounds $\|D^jX_q\|_\infty$.  This specifies constants using only
finite differentiation and the displayed parameters.

The envelope identity is
$D S_q=\partial_z\mathcal A(z,X_q(z))$.  Differentiating it through
order $m+1$ uses only the already bounded derivatives of $X_q$ and
of $L$ through order $m+2$.  Together with the bound for $S_q$ itself
this proves the asserted $C^{m+2}$ estimate.  Composition and the
$C^m$ product inequality now apply to~\eqref{eq:v80-weight-budget}:
the momentum composition uses derivatives of $S_q$ through $m+1$,
and its Jacobian uses derivatives through $m+2$.  The mixed derivative
has a fixed nonzero sign, so taking its absolute value creates no loss
of regularity.  The lower bound is
$D^{-1}e_-\rho_-u_q$.  Multiplication by $T_j-S_q$ and division by
a positive scalar mass prove the remaining assertions.
\end{proof}

\begin{theorem}[Uniform budget and four-deadline convex reconstruction]
\label{thm:v80-four-convex}
Fix $R,n$ and the class constants of~\eqref{eq:v79-convex-class}.
Use the phase box, gates, roof offset $c$, release interval and the
first two deadlines at each of the counts $n,n+1$ constructed in
Theorem~\ref{thm:v79-six-clock}.  The four \emph{raw} endpoint laws
of the unknown contact suspension determine $L|_{[-R,R]^2}$, including
its additive constant.  All settings are independent of $L$.

For a target $C^k$ estimate, $k\ge0$, it suffices to impose fixed
$C^{k+4}$ bounds on $L$ on the enlarged configuration square and
fixed $C^{k+2}$ bounds on the raw source and efficiency on the enlarged
phase and endpoint domains, with positive source and efficiency floors.
Under these explicit bounds the inverse is uniformly locally Lipschitz
from $C^{k+2}$ raw subdensities to $C^k([-R,R]^2)$.
Clock-dependent log-efficiency errors of $C^{k+2}$ size $\epsilon$
add at most $C\epsilon$ to the bound.  Full-rank calibrated branch
channels and their errors can be included as in
Theorem~\ref{thm:v80-channel} whenever the declared finite sheets form
an exhaustive acquisition block.
\end{theorem}
\begin{proof}
The class-wide construction gives true raw densities
$a_q(T_j-S_q)$ for $q=n,n+1$, with the same normalized source at
the two deadlines of each count.  Theorem~\ref{thm:v80-two-raw}
recovers both absolute actions.  The unique physical prefix match of
Lemma~\ref{lem:v79-convex-twist} then gives
\[
       L(t,v)=S_{n+1}(\sigma(t,v),v)-S_n(\sigma(t,v),t)
\]
on the whole prescribed square.  This is the same deterministic inverse
as in the six-conditional-law theorem; success masses have replaced
its third clocks.

Apply Proposition~\ref{prop:v80-stationary-bounds} with $m=k+2$.
It gives exactly the action and raw-density bounds required for the
$C^{k+2}$ inverse and its weight floor.  The matching derivative at a
root has magnitude at least $\mu=u_n^2/(2\Lambda)$.  The resulting
uniform $C^3$ action bound gives a uniform Lipschitz constant $M_3$ for
this derivative in the initial coordinate.  Every true root has the
unit initial buffer supplied by the backward construction.  A collar
of radius at most $\min(1/2,\mu/(2(1+M_3)))$ therefore preserves
$|H_s|\ge\mu/2$; its endpoints have opposite signs separated from
zero by this derivative bound times the collar radius.  These are
uniform root collars, not unspecified higher-derivative assumptions.
The implicit and composition estimate in
Proposition~\ref{prop:v77-cancellation} proves the asserted $C^k$
bound.  The raw nuisance estimate and, when present, the calibrated
channel estimate supply the two extensions.
\end{proof}

For density smoothing with an additional $s$ derivatives, replace the
bounds in this theorem by $C^{k+4+s}$ for $L$ and $C^{k+2+s}$ for
source and efficiency.  These are sufficient budgets; optimality of
the derivative loss is not claimed.  The six-conditional-law result is
retained because a conditional-law archive may no longer contain the
success masses needed by the four-raw-law theorem.
